% !TEX program = xelatex
\documentclass{misisarticle}

\grnti{20.23.17}
\asjc{1710}
\oecd{1.2}

\title{Алгоритм ранжирования конфигураций HTAP-систем по требованиям смешанной нагрузки}

\misisauthor{Д.~С.~Калинкин}{Кафедра инженерной кибернетики НИТУ «МИСИС»}{Москва, Россия}{m2302322@edu.misis.ru}
\misisauthor{А.~В.~Соловьев}{Кафедра инженерной кибернетики НИТУ «МИСИС»}{Москва, Россия}{soloviev@isa.ru}

\abstract{В работе решается задача выбора конфигурации HTAP-системы для смешанной транзакционной и аналитической нагрузки. Предложен алгоритм RASH, который учитывает задержки OLTP- и OLAP-запросов, свежесть данных, стоимость ресурсов и разброс результатов между запусками. Веса критериев вычисляются из параметров нагрузки и заданных ограничений. Сначала алгоритм отделяет конфигурации, которые выполняют все SLO, а затем сравнивает варианты по Парето-слоям и итоговому баллу. Программный прототип проверен на трёх синтетических профилях, микробенчмарке SQLite и DuckDB, а также на смешанной нагрузке BenchBase CH-benCHmark для PostgreSQL. Анализ чувствительности показал, что при изменении требований лучший вариант также меняется.}
\keywords{HTAP, OLTP, OLAP, смешанная нагрузка, многокритериальная оптимизация, ранжирование конфигураций}

\begin{document}
\maketitle

\section{Введение}

	Информационные системы часто должны одновременно выполнять короткие транзакции и строить отчёты по недавно изменённым данным. Если разделить эти задачи между OLTP-системой и отдельным хранилищем, каждую часть проще настраивать, но появляется задержка переноса данных. HTAP-системы стараются совместить транзакционную и аналитическую обработку. Однако требования всё равно конфликтуют: транзакциям нужна малая и стабильная задержка, а аналитическим запросам требуется быстрое чтение больших объёмов данных.

	Современные HTAP-архитектуры различаются способом хранения строк и столбцов, изоляцией вычислительных ресурсов, механизмом репликации и правилами маршрутизации запросов~\cite{zhangSurvey2024}. Поэтому выбор конфигурации нельзя свести к сравнению одной метрики. Ускорение аналитики с помощью отдельной колоночной реплики увеличивает стоимость и может ухудшить свежесть данных. Выполнение аналитики в основном хранилище уменьшает лаг, но усиливает взаимное влияние нагрузок.

	На практике требования обычно задают через SLO: максимальную p99-задержку транзакций, допустимое время аналитического запроса, лаг данных и бюджет. Если просто сложить показатели с разными весами, хорошее значение одной метрики может скрыть нарушение другой. Например, дешёвая, но слишком медленная конфигурация может оказаться выше варианта, который выполняет все обязательные требования.

	Цель работы состоит в разработке воспроизводимого способа ранжирования HTAP-конфигураций по требованиям конкретной нагрузки. В работе предложен алгоритм RASH (Requirement-Aware Stable HTAP Ranking), формат входных измерений и прототип с синтетическим стендом и микробенчмарком. Элемент новизны состоит в том, что в одной процедуре объединены автоматический расчёт весов, обязательная проверка SLO, Парето-ранжирование и учёт разброса повторных запусков.

\section{Анализ работ}

	В обзоре Ч.~Чжана и соавторов HTAP-системы разделены на четыре класса: основное строковое хранилище с колоночным представлением в памяти, распределённое строковое хранилище с колоночной репликой, строковое хранилище с распределённым колоночным контуром и основное колоночное хранилище с дельта-строками~\cite{zhangSurvey2024}. Авторы выделяют обработку смешанной нагрузки, организацию и синхронизацию данных, оптимизацию запросов и планирование ресурсов. Такая классификация показывает, что сравнивать следует конфигурации, а не только названия СУБД.

	Бенчмарк HATtrick оценивает скорость транзакционной и аналитической обработки, а также свежесть данных~\cite{milkai2022}. Он показывает, как два вида нагрузки мешают друг другу. В данной работе используется более простой набор метрик: p99 OLTP, p95 OLAP, лаг свежести и стоимость. RASH не заменяет HATtrick, а может использовать полученные им измерения как входные данные.

	CH-benCHmark объединяет транзакционный профиль TPC-C и аналитические запросы на основе TPC-H. Генератор BenchBase позволяет задавать интенсивность и смесь операций, а также сохраняет задержки разных типов транзакций~\cite{difallah2013,benchbase}. В данной работе BenchBase используется как дополнительная проверка на одновременной нагрузке PostgreSQL. Свежесть данных этот запуск не измеряет, поэтому его результаты рассматриваются отдельно от полного ранжирования RASH.

	Для небольшого реального теста выбраны SQLite и DuckDB. SQLite хранит данные построчно и хорошо подходит для коротких транзакционных операций~\cite{gaffney2022}. DuckDB использует колоночное хранение и векторизованное выполнение, поэтому ориентирован на аналитику~\cite{duckdb2019}. Эти СУБД не являются полноценными распределёнными HTAP-платформами. Однако их различия позволяют на одном компьютере проверить, как алгоритм реагирует на задержки записи, аналитики и синхронизации.

	Автоматическую настройку СУБД часто решают как задачу оптимизации гиперпараметров. С.~Чжан и соавторы выделяют три части этого процесса: выбор важных параметров, поиск конфигурации и перенос результатов между нагрузками~\cite{zhangTuning2022}. Такие методы полезны при большом числе настроек, но требуют многих запусков СУБД. В данной работе задача проще: варианты уже заданы, их нужно понятно упорядочить по требованиям нагрузки.

	Работы О.~Хаммуда, И.~Тарханова и А.~Космарского демонстрируют экспериментальное сравнение вариантов распределённого хранилища по скорости, надёжности и объёму ресурсов~\cite{hammoud2021}. В последующей работе рассмотрена модель со встроенной балансировкой нагрузки~\cite{hammoud2022}. Эти исследования относятся к другой предметной области, однако подтверждают практическую ценность явного сопоставления архитектурных вариантов по нескольким показателям. В методике И.~А.~Тарханова и соавторов один показатель эффективности может быть целевым, а остальные используются как ограничения~\cite{tarkhanov2020}. Этот принцип применён в RASH при разделении допустимости и итогового балла.

	Таким образом, существующие работы предлагают хорошие средства измерения и автоматического поиска настроек. RASH решает более узкую задачу: ранжирует небольшой заранее заданный набор HTAP-конфигураций по конкретным SLO. По сложности этот подход находится между ручным сравнением в таблице и полноценной автоматической оптимизацией.

\section{Постановка задачи}

	Пусть профиль смешанной нагрузки задаётся вектором:
\begin{equation}
 P=(q_t,q_a,\lambda,S_t,S_a,\Delta,C_{\max}),
 \label{eq:profile}
\end{equation}

	где $q_t$ и $q_a$ обозначают доли OLTP- и OLAP-операций, $\lambda$ задаёт суммарную интенсивность поступления операций, $S_t$ задаёт допустимую p99-задержку транзакций, $S_a$ задаёт допустимую p95-задержку аналитических запросов, $\Delta$ задаёт максимальный лаг свежести, а $C_{\max}$ задаёт бюджет за единицу времени. Доли нагрузки нормируются так, что $q_t+q_a=1$.

	Рассматривается множество кандидатов $X=\{x_1,\ldots,x_n\}$. Кандидат описывает архитектурный класс, число процессорных ядер, объём памяти, относительную эффективность OLTP и OLAP, степень изоляции контуров, период синхронизации и стоимость. После $r$ повторов для конфигурации $x_i$ получаем измерения:
\begin{equation}
 M_i^{(j)}=(L_{t,i}^{(j)},L_{a,i}^{(j)},D_i^{(j)},C_i),
 \label{eq:measurement}
\end{equation}

	где $j=1,\ldots,r$, $L_t$ и $L_a$ обозначают наблюдаемые задержки, $D$ обозначает лаг свежести, а $C$ обозначает стоимость.

	Нужно упорядочить все конфигурации от лучшей к худшей. При этом вариант, выполняющий все SLO, должен находиться выше нарушителя. Важность OLTP и OLAP должна зависеть от профиля нагрузки. Кроме того, решение нельзя основывать на одном удачном запуске, а одинаковые входные данные должны давать одинаковый результат.

	Прототип только ранжирует готовые варианты и не создаёт новые настройки СУБД. Список кандидатов задаёт пользователь или другая программа. Это уменьшает число экспериментов и позволяет отдельно проверить сам порядок выбора.

\section{Описание метода}

\subsection{\textnormal{Вычисление весов}}

	Веса критериев рассчитываются из профиля нагрузки, поэтому пользователю не нужно вручную задавать четыре отдельных числа. Сначала вычисляются коэффициенты:
\begin{align}
 a_t &= 0{,}10+0{,}65q_t+0{,}15p,\\
 a_a &= 0{,}10+0{,}65q_a+0{,}10p,\\
 a_d &= 0{,}10+0{,}25q_a+\frac{0{,}30}{1+\Delta/10},\\
 a_c &= 0{,}10+\frac{0{,}30}{1+C_{\max}/10},
 \label{eq:rawweights}
\end{align}

	где $p=\min(\lambda/1000,1)$. Затем коэффициенты нормируются: $w_k=a_k/\sum a_k$. Постоянная $0{,}10$ оставляет ненулевой вес у каждого критерия. Остальные коэффициенты выбраны для первой версии алгоритма. В дальнейшем их следует уточнить по результатам реальных экспериментов или экспертных оценок.

\subsection{\textnormal{Нормирование и проверка SLO}}

	Для средних по повторам метрик вычисляются безразмерные отношения:
\begin{equation}
 r_i=\left(\frac{\overline L_{t,i}}{S_t},
 \frac{\overline L_{a,i}}{S_a},
 \frac{\overline D_i}{\Delta},
 \frac{C_i}{C_{\max}}\right).
 \label{eq:ratios}
\end{equation}

	Если все компоненты $r_i$ не превышают единицу, конфигурация выполняет заданные SLO. Её базовый балл рассчитывается как взвешенное геометрическое среднее:
\begin{equation}
 B_i=\prod_{k=1}^{4}\max(r_{ik},\varepsilon)^{w_k},
 \label{eq:base}
\end{equation}

	где $\varepsilon=10^{-12}$ защищает вычисление от нулевых значений. Геометрическое среднее работает с относительными величинами и не позволяет одному очень хорошему показателю полностью скрыть плохой.

	Для нарушителей используется квадратичный штраф:
\begin{equation}
 P_i=10\sum_{k=1}^{4}w_k\max(0,r_{ik}-1)^2.
 \label{eq:penalty}
\end{equation}

	Пока SLO выполняется, штраф равен нулю. После превышения границы он быстро растёт. Кроме того, при сортировке сначала проверяется сам факт выполнения SLO. Поэтому низкая стоимость не может скрыть даже небольшое нарушение обязательного требования.

\subsection{\textnormal{Учет устойчивости}}

	Для каждого критерия вычисляется коэффициент вариации $cv_{ik}=\sigma_{ik}/\overline m_{ik}$. Сводная нестабильность равна:
\begin{equation}
 V_i=\frac{1}{4}\sum_{k=1}^{4}cv_{ik}.
 \label{eq:variation}
\end{equation}

	Стоимость в текущем прототипе постоянна между повторами и даёт нулевой вклад. Итоговый численный балл, который требуется минимизировать, имеет вид:
\begin{equation}
 Q_i=B_i+0{,}25V_i+P_i.
 \label{eq:score}
\end{equation}

	Поэтому при близких средних значениях выше окажется конфигурация, которая показывает более стабильный результат между запусками.

\subsection{\textnormal{Парето-слои и порядок}}

	Перед итоговой сортировкой строятся Парето-слои. Конфигурация $x_i$ доминирует $x_j$, если она не хуже по всем четырём метрикам и лучше хотя бы по одной. Варианты, над которыми нет доминирующих конфигураций, попадают в первый слой. Затем они убираются, и процедура повторяется. Сначала кандидаты сортируются по выполнению SLO, затем по номеру слоя и баллу $Q_i$. Имя используется только тогда, когда все остальные значения совпали.

	При фиксированном числе критериев вычисление средних требует $O(nr)$ операций, а сложность построения Парето-слоёв составляет $O(n^2)$. Для десятков заранее выбранных конфигураций такая сложность приемлема. При переходе к тысячам вариантов алгоритм целесообразно объединить с предварительным отбором или методом оптимизации гиперпараметров.

\subsection{\textnormal{Алгоритм ранжирования}}

	Вычислительная процедура RASH состоит из следующих шагов:
\begin{enumerate}[leftmargin=*,label=\arabic*)]
	\item проверить полноту профиля, положительность SLO и положительную сумму долей $q_t+q_a$, затем нормировать доли;
	\item для каждого кандидата агрегировать повторы и вычислить отношения~\eqref{eq:ratios};
	\item получить веса по~\eqref{eq:rawweights}, признак допустимости, $B_i$, $V_i$ и $P_i$;
	\item построить Парето-слои по четырём средним метрикам;
	\item отсортировать кортежи $(\neg feasible_i, layer_i,Q_i,name_i)$ по возрастанию.
\end{enumerate}

	Последний шаг отличает RASH от простой сортировки по общему баллу. Например, в сбалансированном профиле распределённая реплика получила $Q=0{,}299$, но стоила 8{,}4 при бюджете 8. Другая конфигурация имела худший балл $Q=0{,}370$, зато выполняла все требования. Сортировка только по $Q$ выбрала бы нарушителя. В RASH допустимый вариант всегда ставится выше.

\subsection{\textnormal{Программная реализация}}

	Прототип написан на Python; ядро ранжирования использует только стандартную библиотеку, а для натурного стенда подключается DuckDB. Команда \texttt{run} загружает профиль и набор кандидатов, проводит заданное число повторов, сохраняет измерения в JSON и формирует CSV с ранжированием. Команда \texttt{rank} повторно обрабатывает ранее сохранённые или импортированные измерения. Генератор случайных чисел инициализируется строкой из seed, имени профиля и кандидата, поэтому результат не зависит от порядка кандидатов.

	В синтетической модели задержки зависят от интенсивности нагрузки, доступных ресурсов и эффективности каждого контура. При слабой изоляции разброс результатов увеличивается. Модель не повторяет работу конкретной СУБД, а нужна для проверки логики алгоритма и формата данных. Дополнительно прототип умеет разбирать исходные файлы BenchBase и выделять p99 OLTP и p95 OLAP.

\section{Вычислительная апробация}

\subsection{\textnormal{Имитационный эксперимент}}

	Для проверки сформированы четыре условные конфигурации: единое строковое хранилище, строковое хранилище с колоночной репликой, гибридное хранилище и распределённая система с колоночной репликой. Каждый вариант запускался по 30 раз с seed 20261031. Использованы три профиля: транзакционный, аналитический и сбалансированный. Числа являются результатом имитационной модели, а не измерениями PostgreSQL, ClickHouse, TiDB или другой реальной СУБД. В таблице~\ref{tab:results} представлены победители для каждого профиля.

\begin{table}[htbp]
  \caption{Победители для трёх профилей нагрузки}
  \label{tab:results}
  \centering
  \small
  \begin{tabular}{@{}llr@{}}
    \toprule
    \textbf{Профиль} & \textbf{Конфигурация} & \textbf{$Q$} \\
    \midrule
    OLTP & row-store & 0{,}640 \\
    OLAP & распределённая реплика & 0{,}274 \\
    Смешанный & гибрид & 0{,}281 \\
    \bottomrule
  \end{tabular}
\end{table}

	В транзакционном профиле вес OLTP равен 0{,}450, а бюджет ограничен 3 условными единицами в час. Строковое хранилище оказалось единственным вариантом, выполнившим все ограничения: p99 OLTP составила 193{,}964~мс при SLO 210~мс, а стоимость составила 2{,}2. Более производительные варианты исключены из допустимой группы из-за бюджета.

	В аналитическом профиле вес OLAP увеличился до 0{,}395, а его SLO был ужесточён до 600~мс. Распределённая реплика показала p95 OLAP 349{,}186~мс и заняла первое место. У гибридного хранилища p95 составила 702{,}864~мс, поэтому низкий лаг свежести не компенсировал нарушение аналитического SLO.

	В сбалансированном профиле победило гибридное хранилище. Его средние значения составили 71{,}487~мс для p99 OLTP, 664{,}889~мс для p95 OLAP и 0{,}349~с для лага. Стоимость 5{,}5 не превысила бюджет 8 единиц. Повторный запуск с тем же seed сформировал побайтно одинаковые JSON- и CSV-файлы.

	Результаты показывают, что одной лучшей конфигурации для всех профилей нет. Но симуляция проверяет в основном логику RASH, поэтому дополнительно был проведён небольшой эксперимент с реальным выполнением SQL.

\subsection{\textnormal{Микробенчмарк SQLite и DuckDB}}

	Стенд сравнивает SQLite, DuckDB и связку SQLite+DuckDB. Для каждого варианта создавалась одинаковая таблица из 50\,000 строк. Транзакционный тест включал 200 точечных обновлений по первичному ключу, а аналитический тест включал семь запросов с фильтрацией и группировкой. Каждый тест повторялся семь раз. Для ранжирования использовались p99 времени обновления, p95 времени аналитического запроса и средняя задержка появления изменений в аналитической части.

	Обе СУБД работали в памяти и в одном процессе; для DuckDB использовался один поток. Это уменьшает влияние диска и позволяет быстро повторить опыт. При этом стенд не учитывает одновременную работу нескольких клиентов и надёжную запись на диск. В связке SQLite+DuckDB изменения переносились построчно простым кодом на Python. Стоимость задавалась условными коэффициентами 1{,}0, 1{,}2 и 2{,}0, а не реальной ценой облачного сервиса.

	Эксперимент выполнен на Intel Core 5 120U, 12 логических процессорах и 15~ГиБ оперативной памяти, с Python 3.14.4, SQLite 3.46.1 и DuckDB 1.5.6. Профиль содержал равные доли OLTP и OLAP; SLO составляли 2~мс для p99 OLTP, 50~мс для p95 OLAP, 100~мс для свежести и 2{,}1 для ресурсоёмкости. В таблице~\ref{tab:live-results} представлены измеренные значения и балл RASH.

\begin{table}[htbp]
	\caption{Результаты натурного микробенчмарка}
	\label{tab:live-results}
	\centering
	\scriptsize
	\begin{tabular}{@{}lrrrr@{}}
		\toprule
		\textbf{Конфигурация} & \textbf{OLTP} & \textbf{OLAP} & \textbf{Лаг} & \textbf{$Q$} \\
		\midrule
		SQLite & 0{,}010 & 26{,}377 & 0{,}035 & 0{,}065 \\
		DuckDB & 1{,}577 & 3{,}716 & 0{,}717 & 0{,}205 \\
		SQLite+DuckDB & 0{,}017 & 2{,}448 & 171{,}274 & 1{,}722 \\
		\bottomrule
		\multicolumn{5}{@{}l@{}}{Примечание: все значения времени указаны в миллисекундах.}
	\end{tabular}
\end{table}

	\newpage
	В этом тесте SQLite выполнил точечные обновления примерно в 164 раза быстрее DuckDB. В то же время аналитический запрос в DuckDB выполнялся в 7{,}1 раза быстрее. Связка двух СУБД показала лучшее время аналитики, но синхронизация заняла 171{,}274~мс и нарушила SLO свежести. Поэтому RASH поставил этот вариант ниже двух допустимых. SQLite занял первое место, потому что его 26{,}377~мс уже укладывались в аналитический SLO 50~мс и дополнительные ресурсы не давали необходимого преимущества.

\subsection{\textnormal{BenchBase CH-benCHmark}}

	Для дополнительной проверки использован BenchBase CH-benCHmark на PostgreSQL 17.6. Масштаб TPC-C равен одному складу. Одновременно работали по одному терминалу TPC-C и CH-benCHmark. Перед каждым измерением выполнялся прогрев в течение 10 секунд, затем нагрузка измерялась 60 секунд. Опыт повторялся три раза на той же загруженной базе. PostgreSQL был запущен в Docker на том же компьютере, а уровень изоляции был установлен в \texttt{TRANSACTION\_SERIALIZABLE}. В таблице~\ref{tab:benchbase-results} представлены результаты трёх повторов и их медиана.

\begin{table}[htbp]
	\caption{Результаты BenchBase CH-benCHmark}
	\label{tab:benchbase-results}
	\centering
	\scriptsize
	\begin{tabular}{@{}lrrr@{}}
		\toprule
		\textbf{Повтор} & \textbf{p99 OLTP, мс} & \textbf{p95 OLAP, мс} & \textbf{Запр./с} \\
		\midrule
		1 & 11{,}574 & 393{,}114 & 204{,}33 \\
		2 & 8{,}959 & 429{,}069 & 207{,}27 \\
		3 & 8{,}873 & 430{,}721 & 206{,}63 \\
		Медиана & 8{,}959 & 429{,}069 & 206{,}63 \\
		\bottomrule
	\end{tabular}
\end{table}

	В каждом повторе BenchBase выполнил около 12 тысяч транзакций TPC-C и от 261 до 453 аналитических запросов. Разброс p99 OLTP между повторами оказался заметнее, чем разброс p95 OLAP. Это подтверждает необходимость повторных запусков и штрафа за нестабильность. Однако CH-benCHmark в данном опыте не измеряет лаг репликации и стоимость, поэтому по этим данным нельзя вычислить полный балл RASH без дополнительных измерений. HATtrick остаётся источником методики измерения свежести и направлением дальнейшей проверки.

\subsection{\textnormal{Анализ чувствительности}}

	Чтобы проверить влияние выбранных границ, те же измерения были обработаны для 20 профилей. SLO аналитики менялся от 3 до 50~мс, а SLO свежести менялся от 50 до 300~мс. При значениях 3 и 100~мс ни один кандидат не выполнял все требования. Если увеличить допустимый лаг до 200~мс, лучшей становилась связка SQLite+DuckDB. При аналитическом SLO от 5 до 20~мс побеждал DuckDB, а при 50~мс побеждал SQLite. Таким образом, лидер меняется по понятной причине: вместе с требованиями меняется допустимый набор конфигураций.

	По небольшому тесту SQLite и DuckDB нельзя делать вывод, что одна СУБД всегда лучше другой. Таблица мала, нет сети, журналирования и длительного прогрева. Перенос строк на Python намного проще промышленных механизмов CDC, а ресурсные коэффициенты не равны реальной стоимости. Запуск BenchBase добавляет конкурентную смешанную нагрузку, но выполнен только на PostgreSQL и при малом масштабе. Оба эксперимента показывают работу прототипа, но не заменяют полный цикл HATtrick с измерением свежести.

\section{Заключение}

	В работе предложен алгоритм RASH для выбора HTAP-конфигурации по профилю нагрузки: сначала он проверяет обязательные SLO, а затем сравнивает варианты по весам критериев, Парето-слоям и стабильности измерений. Прототип выполняет весь процесс от получения данных до сохранения итогового рейтинга. 
	
	В симуляции для трёх профилей победили разные архитектуры. Тест SQLite и DuckDB показал, что лучший вариант меняется вместе с SLO аналитики и свежести. Три запуска BenchBase подтвердили, что прототип может обрабатывать результаты конкурентной OLTP- и OLAP-нагрузки. 
	
	В дальнейшем алгоритм нужно проверить на нескольких серверных СУБД с измерением свежести по методике HATtrick, а его коэффициенты нужно уточнить по результатам этих экспериментов.

\begin{thebibliography}{99}

\bibitem{zhangSurvey2024} C. Zhang, G. Li, J. Zhang, X. Zhang, and J. Feng, ``HTAP Databases: A Survey,'' \textit{IEEE Transactions on Knowledge and Data Engineering}, vol.~36, no.~11, pp.~6410--6429, 2024, doi: \href{https://doi.org/10.1109/TKDE.2024.3389693}{10.1109/TKDE.2024.3389693}.

\bibitem{milkai2022} E. Milkai, Y. Chronis, K. P. Gaffney, Z. Guo, J. M. Patel, and X. Yu, ``How Good is My HTAP System?,'' in \textit{Proceedings of the 2022 International Conference on Management of Data}, pp.~1810--1824, 2022, doi: \href{https://doi.org/10.1145/3514221.3526148}{10.1145/3514221.3526148}.

\bibitem{difallah2013} D. E. Difallah, A. Pavlo, C. Curino, and P. Cudr\'{e}-Mauroux, ``OLTP-Bench: An Extensible Testbed for Benchmarking Relational Databases,'' \textit{Proceedings of the VLDB Endowment}, vol.~7, no.~4, pp.~277--288, 2013, doi: \href{https://doi.org/10.14778/2732240.2732246}{10.14778/2732240.2732246}.

\bibitem{benchbase} Carnegie Mellon Database Group, ``BenchBase: Multi-DBMS SQL benchmarking framework via JDBC,'' GitHub repository. [Online]. Available: \url{https://github.com/cmu-db/benchbase}. Accessed: Nov.~5, 2026.

\bibitem{gaffney2022} K. P. Gaffney, M. Prammer, L. Brasfield, D. R. Hipp, D. Kennedy, and J. M. Patel, ``SQLite: Past, Present, and Future,'' \textit{Proceedings of the VLDB Endowment}, vol.~15, no.~12, pp.~3535--3547, 2022, doi: \href{https://doi.org/10.14778/3554821.3554842}{10.14778/3554821.3554842}.

\bibitem{duckdb2019} M. Raasveldt and H. M\"{u}hleisen, ``DuckDB: An Embeddable Analytical Database,'' in \textit{Proceedings of the 2019 International Conference on Management of Data}, pp.~1981--1984, 2019, doi: \href{https://doi.org/10.1145/3299869.3320212}{10.1145/3299869.3320212}.

\bibitem{zhangTuning2022} X. Zhang, Z. Chang, Y. Li, H. Wu, J. Tan, F. Li, and B. Cui, ``Facilitating Database Tuning with Hyper-Parameter Optimization: A Comprehensive Experimental Evaluation,'' \textit{Proceedings of the VLDB Endowment}, vol.~15, no.~9, pp.~1808--1821, 2022, doi: \href{https://doi.org/10.14778/3538598.3538604}{10.14778/3538598.3538604}.

\bibitem{hammoud2021} O. Hammoud, I. Tarkhanov, and A. Kosmarski, ``An Architecture for Distributed Electronic Documents Storage in Decentralized Blockchain B2B Applications,'' \textit{Computers}, vol.~10, no.~11, art.~142, 2021, doi: \href{https://doi.org/10.3390/computers10110142}{10.3390/computers10110142}.

\bibitem{hammoud2022} O. Hammoud and I. A. Tarkhanov, ``A Novel Blockchain-Integrated Distributed Data Storage Model with Built-in Load Balancing,'' in \textit{2022 IEEE 16th International Conference on Application of Information and Communication Technologies}, pp.~1--6, 2022, doi: \href{https://doi.org/10.1109/AICT55583.2022.10013548}{10.1109/AICT55583.2022.10013548}.

\bibitem{tarkhanov2020} I. A. Tarkhanov, G. P. Akimova, M. A. Pashkin, and A. V. Soloviev, ``Modelling the Methodology to Assess the Effectiveness of Distributed Information Systems,'' \textit{Advances in Science, Technology and Engineering Systems Journal}, vol.~5, no.~1, pp.~86--92, 2020, doi: \href{https://doi.org/10.25046/aj050111}{10.25046/aj050111}.

\end{thebibliography}

\end{document}
